%% Chương 5: Thực nghiệm
\clearpage
\phantomsection
\section*{CHƯƠNG 5. THIẾT LẬP THỰC NGHIỆM}
\addcontentsline{toc}{section}{\numberline{}CHƯƠNG 5. THIẾT LẬP THỰC NGHIỆM}
\setcounter{section}{5}
\setcounter{subsection}{0}
\setcounter{figure}{0}
\setcounter{table}{0}
\setcounter{equation}{0}

\subsection{Môi trường thực nghiệm}

\subsubsection{Cấu hình phần cứng và phần mềm}

Toàn bộ thực nghiệm được thực hiện trên một máy tính duy nhất nhằm đảm bảo
tính nhất quán trong đo lường hiệu năng:

\begin{table}[ht]
\centering
\caption{Môi trường phần cứng và phần mềm thực nghiệm}
\label{tab:env}
\begin{tabular}{ll}
\toprule
Thành phần & Cấu hình / Phiên bản \\
\midrule
Hệ điều hành        & Ubuntu 22.04 LTS (Linux kernel 6.8.0-124-generic) \\
CPU                  & x86-64 (hỗ trợ AES-NI, SSE4.2) \\
Ngôn ngữ C++         & GCC 11.4, chuẩn C++17, tối ưu \texttt{-O2} \\
Python               & 3.12 với numpy, scipy, scikit-image, opencv-python \\
FFmpeg               & 6.x (decode video, đo PSNR/SSIM frame-level) \\
H.264 analyzer       & h264bitstream \texttt{h264\_analyze} \cite{h264bitstream} \\
Khóa mã hoá         & \texttt{testkey} (ASCII 7 ký tự) \\
\bottomrule
\end{tabular}
\end{table}

\textbf{Lý do chọn cờ biên dịch \texttt{-O2}:}
\texttt{-O2} kích hoạt hầu hết các tối ưu hoá của GCC (loop unrolling, inlining,
register allocation) mà không thay đổi hành vi quan sát được (không như \texttt{-O3}
có thể gây floating-point reordering). Đây là cấp tối ưu cân bằng phổ biến nhất
trong production.

\textbf{Lưu ý về AES-NI:} Cài đặt AES trong nghiên cứu này sử dụng OpenSSL
(hoặc tiny-AES-c), vốn tự động sử dụng AES-NI nếu CPU hỗ trợ. Do đó tốc độ
AES trong thực nghiệm phản ánh hiệu năng tối ưu với hardware acceleration.

\subsubsection{Thư viện và công cụ bên thứ ba}

\begin{itemize}
  \item \textbf{h264bitstream \cite{h264bitstream}:} Thư viện C parse H.264 bitstream,
        cung cấp công cụ \texttt{h264\_analyze} để trích xuất thông tin NALU
        (offset, size, type) từ file .h264 Annex B.

  \item \textbf{OpenSSL:} Thư viện mật mã C/C++, cung cấp implementation AES-128-CTR,
        DES-OFB với hardware acceleration tự động.

  \item \textbf{scikit-image (Python):} Tính SSIM, PSNR trên DC bytes 1D.

  \item \textbf{FFmpeg 6.x:} Decode file H.264 thành raw frames để đo chất lượng
        frame (SSIM frame-level, codec compliance check).
\end{itemize}

\subsection{Dataset}

\subsubsection{File thực nghiệm output.h264}

File thực nghiệm duy nhất là \texttt{videos/output.h264} --- file H.264 raw bitstream
định dạng Annex B (byte stream format):

\begin{table}[ht]
\centering
\caption{Thông số file thực nghiệm output.h264}
\label{tab:dataset}
\begin{tabular}{ll}
\toprule
Thông số & Giá trị \\
\midrule
Kích thước file          & 219 MB (229.376.512 bytes) \\
Định dạng                & H.264 Annex B raw bitstream \\
Profile                  & Baseline Profile \\
Tổng NALU (h264\_analyze) & 19.348 \\
Type 5 (IDR/I-frame)      & 155 \\
Type 1 (P/B-frame)        & 18.882 \\
Non-VCL (SPS, PPS, SEI)  & 311 \\
Số frames video           & 60 \\
NALUs mã hoá S1           & 19.037 (Type 1 + Type 5) \\
NALUs mã hoá S2           & 155 (chỉ Type 5) \\
Tỷ lệ I-frame             & 155 / 19.348 = 0.8\% \\
\bottomrule
\end{tabular}
\end{table}

\subsubsection{Đặc điểm phân phối NALU}

File \texttt{output.h264} có tỷ lệ I-frame = 0.8\% (155/19.348), tương ứng
với GOP (Group of Pictures) size $\approx 125$ --- tức là cứ 125 frames có 1 keyframe.
Đây là cấu hình phổ biến cho video streaming: GOP lớn giúp giảm bitrate
nhưng tăng độ trễ khi random seek.

\textbf{Phân phối DC bytes:}
DC bytes trong video tự nhiên có phân phối gần Laplacian (phân phối hai nhánh
quanh 0 sau quantization). Phân phối này ảnh hưởng đến kết quả UACI vì
$|C_1[i] - C_2[i]|$ phụ thuộc vào giá trị DC gốc.

\subsubsection{Lý do chỉ dùng 1 file thực nghiệm}

Nghiên cứu này giới hạn ở 1 file video để:
\begin{enumerate}
  \item Đảm bảo kiểm soát chặt chẽ các biến độc lập (chỉ thay đổi thuật toán và chiến lược).
  \item Tập trung vào so sánh định lượng giữa các thuật toán trong cùng điều kiện.
  \item Phạm vi mở rộng đa-video được đề xuất như hướng nghiên cứu tiếp theo (Chương 7).
\end{enumerate}

Hạn chế này được thừa nhận: kết quả không đại diện cho mọi loại video (resolution,
bitrate, content type). Tuy nhiên, cấu trúc NALU và DC offset không thay đổi theo
nội dung video (chỉ thay đổi theo H.264 profile), nên tính đúng đắn
(BER = 0, SHA-256 PASS) có thể tổng quát hoá.

\subsection{Cấu trúc pipeline C++}

\subsubsection{Tổng quan pipeline}

Pipeline end-to-end được triển khai hoàn toàn bằng C++ (file đọc/ghi, NALU parsing,
mã hoá/giải mã) ngoại trừ bước đo metrics (Python):

\begin{figure}[ht]
\centering
%% TikZ: End-to-end pipeline diagram
\begin{tikzpicture}[
  node distance=0.5cm and 0.3cm,
  step/.style={draw, rectangle, rounded corners=3pt,
               minimum width=2.8cm, minimum height=0.65cm,
               text centered, font=\small, fill=blue!10},
  file/.style={draw, rectangle, rounded corners=2pt,
               minimum width=2.8cm, minimum height=0.6cm,
               text centered, font=\small\ttfamily, fill=gray!12},
  verify/.style={draw, diamond, aspect=2.2,
                 minimum width=2.5cm, minimum height=0.6cm,
                 text centered, font=\small, fill=green!15},
  metrics/.style={draw, rectangle, rounded corners=3pt,
                  minimum width=2.8cm, minimum height=0.65cm,
                  text centered, font=\small, fill=orange!20},
  arrow/.style={->, >=stealth, thick},
  label/.style={font=\footnotesize\itshape, gray},
]

%% Column: Encrypt path
\node[file]   (orig)    {output.h264 (gốc)};
\node[step, below=0.4cm of orig]  (extract) {Extract NALU cache};
\node[step, below=0.4cm of extract] (enc)   {Encrypt DC bytes\\(IEncryption)};
\node[file, below=0.4cm of enc]  (encf)   {output.enc (đã mã hoá)};
\node[step, below=0.4cm of encf]  (dec)    {Decrypt DC bytes\\(IEncryption)};
\node[file, below=0.4cm of dec]  (decf)   {output.dec (giải mã)};
\node[verify, below=0.4cm of decf] (cmp)   {cmp byte-exact?};
\node[metrics, below=0.6cm of cmp] (met)   {measure\_metrics.py\\(security/perf/...)};
\node[file, below=0.4cm of met]  (report) {metrics\_report.txt};

%% Arrows
\draw[arrow] (orig)    -- (extract);
\draw[arrow] (extract) -- (enc) node[midway, right, label] {nalu\_info.bin};
\draw[arrow] (orig)    -- ++(3.2,0) |- (enc);
\draw[arrow] (enc)     -- (encf);
\draw[arrow] (encf)    -- (dec);
\draw[arrow] (orig)    -- ++(3.2,0) |- (dec);
\draw[arrow] (dec)     -- (decf);
\draw[arrow] (decf)    -- (cmp);
\draw[arrow] (orig)    -- ++(3.2,0) |- (cmp);
\draw[arrow] (cmp)     -- node[right, label] {PASS} (met);
\draw[arrow] (encf)    -- ++(3.2,0) |- (met);
\draw[arrow] (decf)    -- ++(3.2,0) |- (met);
\draw[arrow] (orig)    -- ++(3.2,0) |- (met);
\draw[arrow] (met)     -- (report);

%% FAIL branch
\node[right=3.5cm of cmp, font=\small, color=red] (fail) {FAIL → debug};
\draw[arrow, red] (cmp.east) -- node[above, font=\footnotesize] {FAIL} (fail);

\end{tikzpicture}

\caption{Pipeline end-to-end: từ file H.264 gốc qua mã hoá/giải mã đến đo lường metrics}
\label{fig:pipeline}
\end{figure}

\subsubsection{Bước 1: Trích xuất và cache thông tin NALU}

\begin{lstlisting}[language=bash, caption={Trích xuất NALU offsets và types bằng h264\_analyze}]
# Bước 1a: Gọi h264_analyze để parse bitstream
./extract/extract_nalu_from_h264analyze videos/output.h264
# Output: nalu_cache/output.h264.nalu_info.bin (binary cache)
\end{lstlisting}

Công cụ \texttt{extract\_nalu\_from\_h264analyze} gọi \texttt{h264\_analyze} từ
h264bitstream library để parse toàn bộ bitstream và trích xuất:
\begin{itemize}
  \item Offset byte tuyệt đối của mỗi NALU (\texttt{nal\_pos}).
  \item Kích thước NALU (\texttt{nal\_size}).
  \item Loại NALU (\texttt{nal\_unit\_type}).
\end{itemize}

Kết quả được serialize thành file nhị phân \texttt{nalu\_cache/output.h264.nalu\_info.bin}.
Cache này cho phép các pipeline mã hoá/giải mã đọc trực tiếp thông tin NALU
mà không cần parse lại bitstream --- giảm chi phí I/O và đảm bảo tính nhất quán.

\textbf{Lý do dùng cache nhị phân:} Parse toàn bộ 19.348 NALU bằng h264\_analyze
mất $\sim$2 giây. Với 8 pipeline cần chạy, cache giúp tiết kiệm $\sim$16 giây.

\subsubsection{Bước 2: Mã hoá (Encryption)}

\begin{lstlisting}[language=bash, caption={Mã hoá DC bytes với strategy S1 và algo hybrid}]
./pipelines/pipeline_hybrid_s1_h264analyze \
    videos/output.h264 testkey --algo hybrid
# Output: results/output.h264/hybrid_s1/output.h264.s1_hybrid_fixed
# Timing: in ra DC_ONLY_TIME_MS sau khi hoàn thành
\end{lstlisting}

Pipeline mã hoá thực hiện các bước:
\begin{enumerate}
  \item Đọc file gốc \texttt{output.h264} vào bộ nhớ (219 MB).
  \item Đọc cache NALU info (\texttt{nalu\_info.bin}).
  \item Với mỗi NALU có type phù hợp (S1: type 1+5; S2: type 5):
    \begin{enumerate}
      \item Trích xuất DC bytes: \texttt{buffer[nal\_pos + 19 .. nal\_pos + 43]}.
      \item Gọi \texttt{IEncryption::encrypt(dc\_bytes, key, nalu\_index)}.
      \item Ghi ciphertext trở lại vào buffer.
    \end{enumerate}
  \item Ghi toàn bộ buffer (đã sửa DC bytes) ra file output.
  \item In \texttt{DC\_ONLY\_TIME\_MS}: thời gian thuần của bước 3 (không tính I/O).
\end{enumerate}

\textbf{Tham số pipeline:}
\begin{itemize}
  \item \texttt{--algo}: \texttt{hybrid} / \texttt{aes} / \texttt{des} / \texttt{rc4}.
  \item \texttt{key}: chuỗi ASCII (\texttt{testkey} trong thực nghiệm).
  \item Strategy (S1/S2): xác định bởi tên binary (\texttt{pipeline\_*\_s1\_*} hoặc \texttt{\_s2\_*}).
\end{itemize}

\subsubsection{Bước 3: Giải mã (Decryption)}

\begin{lstlisting}[language=bash, caption={Giải mã với cùng key và algo}]
./pipelines/pipeline_hybrid_decrypt_s1_h264analyze \
    results/output.h264/hybrid_s1/output.h264.s1_hybrid_fixed \
    testkey --algo hybrid
# Output: .../output.h264.s1_hybrid_fixed.decrypted
\end{lstlisting}

Pipeline giải mã có cấu trúc tương tự mã hoá, chỉ thay \texttt{encrypt()} bằng
\texttt{decrypt()} và đầu vào là file encrypted thay vì file gốc.

Quan trọng: \texttt{nalu\_index} phải giống hệt lúc mã hoá để PLCM sinh đúng
keystream. Đây được đảm bảo bởi việc đọc cùng file cache và duyệt NALU theo
cùng thứ tự.

\subsubsection{Bước 4: Xác minh byte-exact}

\begin{lstlisting}[language=bash, caption={Xác minh byte-exact bằng cmp}]
cmp --silent videos/output.h264 \
    results/output.h264/hybrid_s1/output.h264.s1_hybrid_fixed.decrypted \
    && echo 'BYTE-EXACT OK' || echo 'MISMATCH'
\end{lstlisting}

Lệnh \texttt{cmp} so sánh từng byte hai file (219 MB). Bất kỳ sai lệch nào
đều dừng ngay và báo lỗi. Đây là kiểm tra nhanh trước khi chạy metrics.

\subsubsection{Bước 5: Đo lường metrics}

\begin{lstlisting}[language=bash, caption={Đo metrics cho hybrid S1 với testkey}]
python3 scripts/measure_metrics.py \
    --original  videos/output.h264 \
    --encrypted results/output.h264/hybrid_s1/output.h264.s1_hybrid_fixed \
    --decrypted results/output.h264/hybrid_s1/output.h264.s1_hybrid_fixed.decrypted \
    --strategy S1 --algo hybrid --key testkey
# Output: metrics_report_S1_<timestamp>.txt
\end{lstlisting}

\subsection{Script đo metrics (measure\_metrics.py)}

\subsubsection{Kiến trúc script}

Script Python \texttt{scripts/measure\_metrics.py} đo 4 nhóm chỉ số:

\begin{enumerate}
  \item \textbf{Nhóm bảo mật:} NPCR, UACI (chỉ hybrid), Entropy Shannon,
        SSIM DC, TSSIM DC.
  \item \textbf{Nhóm chất lượng hình ảnh:} PSNR DC (original vs. encrypted).
  \item \textbf{Nhóm tính đúng đắn:} BER, PSNR (original vs. decrypted),
        SHA-256 hash match.
  \item \textbf{Nhóm hiệu năng:} DC\_ONLY\_TIME\_MS (từ pipeline C++ output).
\end{enumerate}

\subsubsection{Đo NPCR/UACI plaintext sensitivity (chỉ hybrid)}

Đây là bước đặc biệt chỉ áp dụng cho hybrid. Script Python tái cài đặt thuật toán
hybrid (bao gồm PLCM, Arnold, và chained diffusion) để sinh hai ciphertext:

\begin{enumerate}
  \item Đọc DC bytes gốc $P$ từ file H.264.
  \item Tạo $P' = P$ với \texttt{byte[0] XOR 1} (flip bit 0 của byte đầu tiên).
  \item Với mỗi NALU: mã hoá cả $P$ và $P'$ với cùng key và \texttt{nalu\_index}.
  \item Tính $D[i]$, $\text{NPCR}$, và $\text{UACI}$ theo công thức 2.3--2.4.
  \item Trung bình NPCR/UACI qua tất cả NALU mã hoá.
\end{enumerate}

Lý do dùng Python thay vì gọi C++ binary: tránh overhead I/O của việc ghi/đọc
file 219 MB hai lần. Thay vào đó, tính toán in-memory trực tiếp.

\textbf{Kiểm tra tính nhất quán:} PLCM Python và PLCM C++ phải cho cùng kết quả.
Điều này được xác minh bằng cách so sánh 10 bytes keystream đầu tiên giữa
Python script và C++ binary cho \texttt{testkey} và \texttt{nalu\_index} = 0.

\subsubsection{Đo Entropy}

\begin{lstlisting}[language=bash, caption={Tính Entropy Shannon từ DC bytes encrypted}]
# Trong measure_metrics.py:
def compute_entropy(data: bytes) -> float:
    counts = np.bincount(np.frombuffer(data, dtype=np.uint8), minlength=256)
    probs = counts[counts > 0] / len(data)
    return float(-np.sum(probs * np.log2(probs)))
\end{lstlisting}

Entropy được tính trên tất cả DC bytes encrypted từ tất cả NALU (không phải per-NALU),
tức là trên mảng khoảng $19.037 \times 25 = 475.925$ bytes (S1) hoặc
$155 \times 25 = 3.875$ bytes (S2).

\subsubsection{Đo SSIM và TSSIM DC}

SSIM DC được tính giữa DC bytes gốc và DC bytes encrypted, coi mỗi NALU là
một tín hiệu 1D 25 phần tử. Trung bình SSIM qua tất cả NALU:

\begin{equation}
\overline{\text{SSIM}}_{\text{DC}} = \frac{1}{N_{\text{enc}}} \sum_{n=1}^{N_{\text{enc}}} \text{SSIM}(\text{DC}_n, C_n)
\end{equation}

TSSIM DC (Temporal SSIM) đo SSIM giữa DC bytes của frame $t$ và frame $t+1$:

\begin{equation}
\overline{\text{TSSIM}}_{\text{DC}} = \frac{1}{N_{\text{enc}}-1} \sum_{t=1}^{N_{\text{enc}}-1} \text{SSIM}(\text{DC}_t^{\text{enc}}, \text{DC}_{t+1}^{\text{enc}})
\end{equation}

\subsubsection{Đo hiệu năng (Timing)}

Thời gian DC encrypt/decrypt được đo từ \textbf{bên trong C++ pipeline} bằng
\texttt{std::chrono::high\_resolution\_clock}:

\begin{lstlisting}[language=C++, caption={Đo thời gian DC encrypt trong pipeline C++}]
auto t_dc_start = std::chrono::high_resolution_clock::now();

// Vòng lặp mã hoá tất cả NALU
for (int i = 0; i < nalu_count; ++i) {
    if (shouldEncrypt(nalu_type[i], strategy)) {
        auto dc = extractDC(buffer, nalu_pos[i]);
        auto enc = cipher->encrypt(dc, key, i);  // IEncryption::encrypt
        writeDC(buffer, nalu_pos[i], enc);
    }
}

auto t_dc_end = std::chrono::high_resolution_clock::now();
double dc_ms = std::chrono::duration<double, std::milli>(t_dc_end - t_dc_start).count();
std::cout << "DC_ONLY_TIME_MS=" << dc_ms << std::endl;
\end{lstlisting}

\texttt{DC\_ONLY\_TIME\_MS} đo \textbf{thuần} thời gian gọi \texttt{encrypt()},
\textbf{không tính} I/O đọc/ghi file 219 MB. Đây là chỉ số so sánh thuật toán
chính xác nhất.

Pipeline time (tổng) = DC time + I/O time + NALU scan time.
I/O 219 MB chiếm phần lớn pipeline time, đặc biệt với S2 (chỉ 155 NALU mã hoá).

\subsection{Quy trình thực nghiệm đầy đủ}

\subsubsection{Chạy tất cả 8 tổ hợp}

\begin{lstlisting}[language=bash, caption={Script đầy đủ chạy 8 tổ hợp thực nghiệm}]
# Bước 1: Build tất cả binaries
make all

# Bước 2: Extract NALU cache (1 lần)
./extract/extract_nalu_from_h264analyze videos/output.h264

# Bước 3: Chạy tất cả 8 tổ hợp encrypt+decrypt+verify
for ALGO in hybrid aes des rc4; do
  for STRAT in s1 s2; do
    # Encrypt
    ./pipelines/pipeline_hybrid_${STRAT}_h264analyze \
        videos/output.h264 testkey --algo $ALGO

    # Decrypt
    ./pipelines/pipeline_hybrid_decrypt_${STRAT}_h264analyze \
        results/output.h264/${ALGO}_${STRAT}/output.h264.${STRAT}_${ALGO}* \
        testkey --algo $ALGO

    # Verify byte-exact
    cmp --silent videos/output.h264 \
        results/output.h264/${ALGO}_${STRAT}/*.decrypted \
        && echo "$ALGO $STRAT: OK" || echo "$ALGO $STRAT: FAIL"
  done
done

# Bước 4: Đo metrics cho từng tổ hợp
for ALGO in hybrid aes des rc4; do
  for STRAT in S1 S2; do
    python3 scripts/measure_metrics.py \
        --original  videos/output.h264 \
        --encrypted results/... \
        --decrypted results/...decrypted \
        --strategy $STRAT --algo $ALGO --key testkey
  done
done
\end{lstlisting}

\subsubsection{Số lần chạy và độ tin cậy}

Mỗi tổ hợp được chạy một lần duy nhất (single run) do thời gian thực nghiệm
lớn (mỗi pipeline $\sim$1--9 giây $\times$ 8 tổ hợp $\times$ 2 (encrypt+decrypt) = tổng $\sim$1 phút).

Timing results (DC\_ONLY\_TIME\_MS) có thể biến động $\pm$5\% do OS scheduling.
Để báo cáo, chúng tôi dùng kết quả của lần chạy cuối cùng sau warm-up (cache đã warm).

\subsection{Tóm tắt thực nghiệm}

\begin{table}[ht]
\centering
\caption{Tóm tắt cấu hình thực nghiệm}
\label{tab:exp_summary}
\begin{tabular}{ll}
\toprule
Tham số & Giá trị \\
\midrule
Số tổ hợp (algo $\times$ strategy) & 8 (4 $\times$ 2) \\
File test & output.h264 (219 MB, 19.348 NALU) \\
Key & \texttt{testkey} \\
Chỉ số đo & 10 (NPCR, UACI, Entropy, SSIM, TSSIM, PSNR, BER, SHA256, DC-time, Pipeline-time) \\
Ngôn ngữ triển khai & C++17 (pipeline) + Python 3.12 (metrics) \\
Thư viện mật mã & OpenSSL (AES, DES), custom (PLCM, Arnold, RC4) \\
\bottomrule
\end{tabular}
\end{table}
